\section{What can be removed from the main paper?}
\label{sec:two-stage-reduction}

The screen-or-solve theorem supports a much shorter primary narrative and
also prevents an unnecessary second stage from becoming the paper's claimed
algorithmic contribution.

\subsection{Proposed main-paper spine}

\begin{enumerate}[leftmargin=2.2em]
\item Point-source PPR, RPPR, semantic error, and degree work.
\item The set-only, numerical, and mass-completion certificate types and
  \cref{alg:two-stage-hybrid}.
\item The numerical early-return lemma, direct APPR/priority-SOR baselines,
  and the evidence that a forced tail is redundant at equal accuracy.
\item The two composition theorems: a fixed linear tail after a set-only
  screen, and APPR/SOR cleanup after a mass-capturing AESP publication.
\item The one-shot hard-capped Green--CG theorem and its fair race with
  direct APPR; then one optional structural theorem (tree or cactus).
\item The envelope-linearization lemma and optional fixed principal PPR solve.
\item The fair certificate race and experiments.
\item The randomized threshold-batch closure, followed by one precise problem
  for deterministic arithmetic and persistent response.
\end{enumerate}

This spine should fit roughly 12--16 pages including proofs and experiments.
The long AESP-CD note remains a proof archive and companion document.

\subsection{Material that moves out of the main line}

The following results remain valuable but are not required to prove the
two-stage composition:
\begin{itemize}[leftmargin=2.2em]
\item correction-by-correction Lyapunov ledgers while momentum crosses faces;
\item Moreau face-shock convolution across many admissions;
\item constant-mode/Perron alignment and rejected-window accounting;
\item mixed-clipping spectral counterexamples;
\item permanent finite-inner momentum replay after a changing-face prefix;
\item high-rank variable-port reporter obstructions not used by the selected
  Stage~I theorem.
\end{itemize}
These belong in appendices or a companion proof-audit paper.  Removing them
from the primary theorem does not discard the research; it prevents an open
changing-face acceleration route from obscuring the closed fixed-envelope
result.

\subsection{The formerly missing and now stronger reporter target}

\begin{problem}[Output-sensitive point-source set-only screening]
\label{prob:two-stage-general-discovery}
On every finite unweighted graph given by adjacency lists, construct an
exact-face or containing-envelope certificate for point-source RPPR using
\begin{equation}
 \widetilde O\!\left(\frac1{\rho\sqrt\alpha}\right)
 \label{eq:two-stage-general-discovery-target}
\end{equation}
total charged work and $\widetilde O(1/\rho)$ retained volume, without global
preprocessing, a KKT margin, or a supplied final support.  The certificate
must be available before the work has already produced a numerical
$\rho$-accurate PPR output; otherwise the correct algorithm is to return that
output and omit Stage~II.
\end{problem}

This exact-or-containing formulation remains a useful stronger set-only
target, but it is no longer necessary for the end-to-end PPR theorem.
Section~\ref{sec:two-stage-randomized-op2} imports the companion randomized
OP2 solver and proves that a certified low-energy \emph{inner} face suffices.
That relaxation closes the arbitrary-graph exact-real randomized PPR bound.
What remains of \cref{prob:two-stage-general-discovery} is its stricter exact
support/containing-set contract, and deterministic finite-precision or
persistent implementations.

By \cref{thm:two-stage-composition}, solving
\cref{prob:two-stage-general-discovery} is sufficient, but no longer
necessary, for the complete point-source PPR target.  The qualification in
the problem is essential:
APPR, priority SOR, gap-certified AESP, and exact structural response all
currently return enough values to invoke
\cref{cor:two-stage-numerical-handoff-returns}.  For those engines the fixed
tail is not a speedup.  The numerical tail and changing-face restart interface
are therefore no longer part of the blocker.  Two genuine reasons for a
continuation remain: a cheaper set-only screen followed by a fixed linear
solve, or an AESP publication that is not yet target-accurate but has only
$O(\sqrt\alpha)$ residual mass and can be completed monotonically.

The most plausible next strategies are:
\begin{enumerate}[leftmargin=2.2em]
\item grow a hard-capped local envelope until the observable Stieltjes
  publication has mass deficit $O(\sqrt\alpha)$, using
  \cref{thm:two-stage-mass-capturing-envelope} instead of identifying
  $S^\star$;
\item a dynamic Schur sign/screening oracle that avoids reconstructing a
  target-accurate numerical obstacle point;
\item a refinement of the hard-capped Green-ball lane that prunes
  low-influence branches without first solving the full numerical target;
\item batched AESP-CD obstacle solves with a clocked gray-row refresh bound;
\item a low-degree response hierarchy combined with high-degree batched
  pivots;
\item homotopy in $\rho$ with persistent response reuse, while charging every
  reorder and boundary notification;
\item a certified portfolio whose numerical lanes return early and whose
  speculative set-only lanes fall back before exceeding an $O(1/\rho)$ hard
  cap.
\end{enumerate}

The mass-capturing route is certificate-wise weaker than support discovery:
it asks
only for a local set whose \emph{principal PPR mass} is large.  It is also
better aligned with AESP-CD, because signed fixed-envelope scratch and safe
publication are already proved.  The remaining graph question is whether an
envelope satisfying that scalar condition can be exposed cheaply on a given
instance.  It cannot be promised at output-sensitive volume on every graph:
\cref{prop:two-stage-star-mass-barrier} makes the diffuse-star obstruction
exact.

Blind radius expansion, repeated full restricted solves after singleton
admissions, and fixed-factor dyadic key refresh are already known to fail as
universal work proofs.  They should not be reintroduced under new names.
